\documentclass[oneside,final,12pt,a4paper]{extreport}
\renewcommand{\thetable}{\Roman{table}}\usepackage{hyperref}
\usepackage{hyperref}

\hypersetup{
    colorlinks=true,
    linkcolor=blue,
    filecolor=magenta,      
    urlcolor=cyan,
    citecolor=cyan,
    pdftitle={Overleaf Example},
    pdfpagemode=FullScreen,
    }
\usepackage[utf8]{inputenc}
\usepackage{graphicx}
\usepackage[affil-it]{authblk}
\usepackage[english]{babel}
\usepackage[backend=biber,style=ieee,autocite=inline]{biblatex}
\usepackage{blindtext}

\usepackage{float}
\usepackage{amsmath,amsfonts}
\usepackage{url}
\usepackage{titlesec}
\usepackage{csquotes}
\usepackage{longtable}
\titleformat{\section}[hang]{\fontsize{20}{24}\selectfont\filcenter}{\Roman{section}}{1em}{}
\titleformat{\subsection}[hang]{\itshape}{\Alph{subsection}.}{1em}{}[]
\titleformat{\subsubsection}[runin]{\itshape}{\arabic{subsubsection})}{1em}{}[$:$]
\titlespacing{\subsubsection}{1em}{1em}{1em}
\titleformat{\paragraph}[runin]{\itshape}{\alph{paragraph})}{1em}{}[$:$\quad]
\titlespacing{\paragraph}{2em}{1em}{1em}
\usepackage{vmargin}
\setpapersize{A4}
\setmarginsrb{2.5cm}{2cm}{2cm}{2cm}{0pt}{10mm}{0pt}{13mm}
\usepackage{setspace}
\sloppy
\setstretch{1.5}
\usepackage{indentfirst}
\usepackage{pgfplots}
\parindent=1.25cm
\usepackage{caption}
\usepackage{subcaption}
\DeclareCaptionLabelSeparator{figSep}{.\quad}
\title{CRISP-DM Report: BiznesDengi}
\author{Aleliya Turushkina, Arina Petuhova, Maksim Ilin, Ivan Ilyichev, Lev Permiakov}


\begin{document}

\maketitle

\newpage

\section{Business Understanding}
\subsection{Determine Business Objectives}


\subsubsection{Background}
The project is initiated by an investment advisory business unit that is responsible for supporting retail clients in making stock market decisions. This unit defines business objectives, evaluates the usefulness of the results, and ensures alignment between analytical outcomes and real-world investment needs. Within the project, it acts as the business customer.

The problem area addressed in this project is short-term investment decision support. Retail investors often lack data-driven tools for selecting stocks that balance return and risk within a limited investment horizon. Existing approaches are typically based on manual analysis or simple heuristics, which may not fully exploit historical market data and can lead to suboptimal investment decisions.

The primary target user group consists of private investors with a six-month investment horizon, aiming to preserve and moderately increase their savings for future expenses. Secondary users include business analysts and investment specialists, who interpret model outputs and translate them into actionable recommendations.

From a business perspective, users expect the project to identify stocks with growth potential, provide interpretable and transparent reasoning behind selections, and support risk-aware investment strategies. In particular, given the short investment horizon, the solution is expected to help manage downside risk while achieving stable returns.

\subsubsection{Business Objective}
The objective of this project is to support short-term investment decision-making by identifying a portfolio of stocks that can provide positive returns over a six-month investment horizon. From a business perspective, the goal is to achieve a stable increase in the investor’s deposit while maintaining an acceptable level of risk.

To make the objective measurable, the target return is set to at least \textbf{11.2\%} over the investment period. This value is motivated by the need to outperform inflation by a significant margin: the official inflation rate in Russia was 5.6\% in 2025 (\href{https://www.cbr.ru/eng/press/event/?id=28252#:~:text=Annual%20inflation%20equalled%205.6%25%20in,rate%20of%203%25%20on%20average.}{Bank of Russia}), and the project aims to achieve at least twice this level to ensure real growth of purchasing power.

At the same time, risk management constraints are introduced: the maximum acceptable loss for an investment strategy is limited to \textbf{20\%} of the initial deposit. This threshold reflects a moderate risk tolerance for retail investors and is supported by observed drawdowns in comparable investment strategies. For example, the “Rational Profile” strategy demonstrates a historical drawdown of approximately 23.38\% (\href{https://gx2invest.ru/strategies/rational}{investment strategy description}), which justifies setting a conservative threshold slightly above this level.

In addition, the solution is expected to provide interpretable recommendations that can be used by analysts and understood by end users. From a business value perspective, achieving these objectives allows investors to outperform inflation and traditional bank deposits, thereby increasing the attractiveness of stock market investments.

\subsubsection{Business Success Criteria}
The project will be considered successful if the proposed investment strategy achieves a return of at least \textbf{11.2\%} over the six-month horizon while not exceeding the maximum acceptable loss of \textbf{20\%}, as defined in the business objective.


\subsection{Situation Assessment}

\subsubsection{Inventory of Resources}
\begin{itemize}
    \item People: We have project manager, business unit, business analyst, data scientist, and machine learning engineer people in the team, so the team is complete for rigorous study.
    \item Data: we will parse 249 most popular Russian stocks, which are easily available from the Moscow Exchange Information and Statistical Service (MOEX ISS) API with different time frames ranging from minutes to months.
    \item Computing resources: Colab and Kaggle are great choices for programming and modeling.
    \item Software: Python with its numerous libraries and possibly Docker for deploying models.
\end{itemize}

\subsubsection{Requirements, Assumptions, and Constraints}
\begin{itemize}
    \item Requirements: historical stocks data for 16 years (15 years for modeling + 1 for evaluating real case scenario). The project duration is estimated to 2 months. Different time frames from minutes to months are needed for in-depth analysis of stocks behavior. The issue is that the Moscow Exchange requires people that are interested to profit from their data, need to contact the officials for negotiating commercial data use (\href{https://www.moex.com/a7156#:~:text=4.3.,by%20the%20Exchange%20and%20Users.}{User Agreement}) (\href{https://www.moex.com/a2193#:~:text=%D0%9F%D1%80%D0%BE%D0%B3%D1%80%D0%B0%D0%BC%D0%BC%D0%BD%D1%8B%D0%B9%20%D0%B8%D0%BD%D1%82%D0%B5%D1%80%D1%84%D0%B5%D0%B9%D1%81%20%D0%BA%20%D0%98%D0%A1%D0%A1,%D1%81%D0%B5%D1%80%D0%B2%D0%B5%D1%80%D0%BE%D0%BC%20%D0%BE%D1%81%D1%83%D1%89%D0%B5%D1%81%D1%82%D0%B2%D0%BB%D1%8F%D0%B5%D1%82%D1%81%D1%8F%20%D0%BF%D0%BE%20%D0%BF%D1%80%D0%BE%D1%82%D0%BE%D0%BA%D0%BE%D0%BB%D1%83%20http.&text=%D0%9D%D0%B0%20Python%20%D0%B8%20Visual%20Basic,NET.}{In Russian}).
    \item Assumptions: the main assumption is that chosen historical period is enough to model profitable investment strategy. Hidden assumption could be that no serious crisis will happen during evaluation year, however this is addressed by that fact that some crises are already included in the data.
    \item Constraints: The project is subject to several practical constraints. First, the timeline is limited to approximately two months, as defined in the project plan, which restricts the complexity of data processing, feature engineering, and model development. 

    Budget is also a constraint, since the project requires human resources, computing equipment, data access, and future promotion costs. Therefore, the solution should remain economically justified relative to the expected business value.
    
    Legal constraints include restrictions on the use of MOEX data, which may require compliance with licensing agreements and exchange data usage policies. In addition, if the system is later applied to real clients, personal data processing and user profiling must comply with privacy regulations such as GDPR or similar data protection standards.
    
    Technical constraints arise from limited computational resources, restricted project duration, and the need to process large volumes of historical financial data within reasonable time limits.
\end{itemize}

\subsubsection{Risks and Contingencies}

Several risks may affect the successful completion and outcomes of the project.

\textbf{Data risks:} The dataset may contain missing values, inconsistencies, or biases (e.g., survivorship bias or incomplete historical records). These issues can negatively impact model training and evaluation. To mitigate this risk, data quality analysis will be performed during the data understanding phase, and appropriate preprocessing techniques such as imputation, filtering, and validation against external sources will be applied.

\textbf{Technical risks:} The developed models may fail to achieve the desired predictive performance or may overfit historical data, leading to poor generalization on unseen data. To address this, multiple modeling approaches will be tested, cross-validation techniques will be applied, and model performance will be continuously evaluated against defined success criteria.

\textbf{Business risks:} The resulting investment strategy may not be adopted by end users due to lack of trust, insufficient interpretability, or misalignment with user expectations. To reduce this risk, emphasis will be placed on model interpretability and clear presentation of results, ensuring that recommendations are understandable and actionable for both analysts and investors.

\textbf{External risks:} Unexpected market conditions, such as financial crises or extreme volatility, may invalidate patterns learned from historical data. As a contingency, risk management rules (e.g., stop-loss constraints) are incorporated, and the limitations of the model will be explicitly communicated.

\subsubsection{Terminology}
\begin{itemize}
    \item Stock - partial ownership in an organization.
    \item Dividend - a sum of money paid by a company to its shareholders out of its profits.
\end{itemize}

\subsubsection{Costs and Benefits}
\textbf{Costs:} The economic feasibility of the project can be assessed by comparing the total implementation and promotion costs with the expected business benefits.

The project cost consists of several components. First, labor costs are estimated using an average junior IT specialist salary of \$1000 per month. For a team of five members working for two months, the total labor cost is
\[
5 \times 1000 \times 2 = \$10{,}000.
\]

Second, infrastructure costs are estimated through laptop depreciation. Assuming that each team member uses a laptop costing \$2000 with a renewal period of four years, the monthly cost per laptop is
\[
2000 / 48 \approx \$41.7.
\]
For five laptops over two months, the total infrastructure cost is
\[
41.7 \times 5 \times 2 \approx \$417.
\]

Third, data acquisition costs are associated with MOEX market data access. Assuming an average cost of 20{,}000 RUB per month and an exchange rate of 1 USD = 80 RUB, the monthly cost is
\[
20000 / 80 = \$250.
\]
For two months, this gives
\[
250 \times 2 = \$500.
\]

Fourth, operational costs such as electricity, internet access, and auxiliary software usage are estimated at approximately \$100 per month for the whole team, resulting in
\[
100 \times 2 = \$200.
\]

In addition, promotion costs should be considered for the future deployment of the solution. Assuming a marketing campaign of 5 paid posts per day in different channels, with an average cost of \$10 per post, the daily marketing cost is
\[
5 \times 10 = \$50.
\]
For three months of promotion (approximately 90 days), the total marketing cost is
\[
50 \times 90 = \$4{,}500.
\]

Thus, the total estimated cost of the project and initial promotion stage is
\[
10{,}000 + 417 + 500 + 200 + 4{,}500 = \$15{,}617.
\]

\textbf{Benefits:} The expected benefits of the project arise from attracting investors to a strategy that aims to outperform inflation and traditional bank deposits. If the solution is adopted by clients, it can generate business value through increased customer engagement, higher trust in investment recommendations, and potential revenue from advisory or portfolio management services. Therefore, the project becomes economically justified if the resulting strategy attracts enough users or managed capital to compensate for the estimated total cost of \$15{,}617.

The expected benefits of the project can be quantified by modeling potential revenue from providing investment recommendations. Let us assume that the business model is based on charging a commission of \textbf{5\%} from the positive return generated for each client.

Assuming an average client investment of \$5,000 and a target return of 11.2\%, the expected profit per client over the six-month horizon is
\[
5000 \times 0.112 = \$560.
\]
The corresponding revenue for the business (5\% commission) is
\[
560 \times 0.05 = \$28 \text{ per client}.
\]

To cover the total project and promotion cost of \$15,617, the required number of clients is
\[
15617 / 28 \approx 558 \text{ clients}.
\]

Thus, the project reaches the break-even point when approximately \textbf{560 clients} adopt the strategy.

If the number of clients exceeds this threshold, the project becomes profitable. For example, with 1000 clients, the expected revenue is
\[
1000 \times 28 = \$28{,}000,
\]
resulting in a net profit of
\[
28000 - 15617 = \$12{,}383.
\]

In terms of payback period, assuming a steady inflow of clients and that each client completes one investment cycle every six months, the project is expected to recover its initial costs within the first cycle if at least 560 clients are acquired. Otherwise, the payback period extends proportionally (e.g., approximately one year if only half of the required clients are reached in the first cycle).

Therefore, the project becomes economically viable if it can attract a moderate number of users, demonstrating that even a relatively small client base can compensate for the initial investment and generate profit.


\subsection{Determine Data Mining Goals}

\subsubsection{Data Mining Goals}

The business objective of achieving stable positive returns with controlled risk is translated into a data mining task focused on predicting stock performance over a six-month horizon.

The primary data mining problem is defined as a \textbf{regression task}, where the goal is to predict the future return of a stock based on historical price data and relevant features. This can be formulated as a time series forecasting problem, where past observations are used to estimate future price movements.

The output of the data mining process is a model that assigns a predicted return (or probability of profitability) to each stock, allowing the construction of an investment portfolio that aligns with the business objective.

\subsubsection{Data Mining Success Criteria}

The performance of the model will be evaluated using technical metrics appropriate for the defined problem types. For the regression task, the primary evaluation metric is \textbf{Mean Absolute Error (MAE)}, which measures the average deviation between predicted and actual returns. Lower MAE values indicate more accurate predictions.

To ensure practical usability, the model must also satisfy interpretability requirements, allowing analysts to understand the key factors influencing predictions. This is important to support trust and adoption of the model in real-world decision-making.


\subsection{Produce Project Plan}

\subsubsection{Project Plan}
\begin{table}[h!]
\centering
\small
\begin{tabular}{|p{2.3cm}|p{1.6cm}|p{2.5cm}|p{2.5cm}|p{2.5cm}|p{2.5cm}|}
\hline
\textbf{Phase} & \textbf{Duration} & \textbf{Resources} & \textbf{Inputs} & \textbf{Outputs} & \textbf{Dependencies} \\
\hline
Business Understanding & 0.5 Weeks & Business Analyst, Project Manager, Business Unit & Business Requirements, Domain Knowledge & Business Objectives, Situation Assessment, Success Criteria, Project Plan & None \\
\hline
Data Understanding & 1.5 Weeks & Data Scientist, Business Analyst, Project Manager & MOEX Stocks Parsed Data & Initial Data Collection Report, Data Description Report, Data Exploration Report, Data Quality Report & Business Understanding \\
\hline
Data Preparation & 1.5 Weeks & Data Scientist, Business Analyst, Project Manager & Raw Data with Description, Exploration, Quality Descriptions & Data Set(s) with Description & Data Understanding \\
\hline
Modeling (Iteration 1) & 1 Week & ML Engineer, Project Manager & Prepared Dataset(s) & Baseline Models, Initial Predictions & Data Preparation \\
\hline
Evaluation (Iteration 1) & 1 Week & Data Scientist, Business Analyst, Business Unit, Project Manager & Baseline Model(s), Test Data & Model Performance Metrics, Business Objectives Alignment Report & Modeling (Iteration 1) \\
\hline
Modeling (Iteration 2) & 1 Week & ML Engineer, Project Manager & Evaluation Feedback, Refined Data & Improved Model(s), Hyperparameter Tuning Results & Evaluation (Iteration 1) \\
\hline
Evaluation (Iteration 2) & 1 Week & Data Scientist, Business Analyst, Business Unit, Project Manager & Improved Model(s) & Final Model Selection, Deployment Readiness Report & Modeling (Iteration 2) \\
\hline
Deployment & 0.5 Weeks & Data Scientist, ML Engineer, Business Analyst & Final Model, Infrastructure Tools & Production System, Monitoring Dashboard and Logs & Evaluation (Iteration 2) \\
\hline
\end{tabular}
\caption{Project Plan}
\label{tab:tab1}
\end{table}

Estimated total project duration is 7-8 weeks (note that 8th week, which is dedicated to deployment, is under discussion) as shown in Table \ref{tab:tab1}

2 modeling iterations are planned as an optimal setup to try twice and analyze what went well and what did not.

\paragraph{Risk Analysis and Schedule Dependencies}

\begin{table}[h!]
\centering
\small
\begin{tabular}{|p{5cm}|p{5cm}|p{5cm}|}
\hline
\textbf{Risk} & \textbf{Action} & \textbf{Schedule Reflection} \\
\hline
Serious Data Quality Issues (Too Many Missing Values and/or Anomalies) & Prepare Data Imputation/Dropping Strategies  & Extends Data Preparation for 1 week  \\
\hline
Insufficient Expressiveness of Features & Collect More Data from Other Sources & Extends Modeling for 1 week and possibly Data Preparation for 1 week (Backtrack to Previous Stages if Needed) \\
\hline
Key People Unavailability & Documentation and involving specialists with similar background & Depends on situation, timeline extension for several weeks \\
\hline
Global Market Changes & Be Flexible and not Reactive to Changes, Understand Shifts First & Freezing the Project or Serious Timeline Extension \\
\hline
\end{tabular}
\caption{Risk Analysis with Schedule Impact and Actions}
\label{tab:tab2}
\end{table}

\par The timeline is presented in the Project Plan Table \ref{tab:tab1}. Data Preparation is the most risky phase for delays due to typical data quality and transformations challenges. This and other risks are detailed in the Table \ref{tab:tab2}

\paragraph{Review Points}
\par Review points are scheduled as follows:
\begin{itemize}
    \item \textbf{Review 1 (Week 0.5):} Business Understanding completion - verify match of data mining goals with business objectives before any work begins.
    \item \textbf{Review 2 (Week 2):} Data Understanding completion - assess data quality. Decide whether to proceed or find additional data.
    \item \textbf{Review 3 (Week 3.5):} Data Preparation completion - validate prepared dataset(s).
    \item \textbf{Review 4 (Week 5.5):} Iteration 1 Evaluation - compare initial business objectives with the results. Go to Business Understanding if mismatch is found.
    \item \textbf{Review 5 (Week 7.5):} Iteration 2 Evaluation - final model(s) selection and preparing for deployment.
    \item \textbf{Review 6 (Week 8):} Deployment completion - monitoring model(s) behavior in real world. 
\end{itemize}


\subsubsection{Initial Assessment of Tools and Techniques}
\par Python has a plenty of great libraries with maintenance, so problems with data processing or modeling tools absence are not likely. Docker is a good deployment software as practice shows.
    \par Trial of simple models first (time series decomposition) seems relatively efficient and explainable from business perspective, and neural networks may enhance the quality of the final result.

\section{Data Understanding}
\subsection{Data Collection}

The dataset was collected by parsing historical market data from the Moscow Exchange Information and Statistical Service (MOEX ISS) API. The data collection was implemented using the \texttt{apimoex} Python library, which is convenient for accessing MOEX market data endpoints. The collected data covers the period from \texttt{01.02.2010} to \texttt{01.02.2026} and includes daily trading records, historical dividend payments, candlestick data across multiple time intervals, and securities description for the predefined list of \textbf{249} most popular Russian stocks.

\subsection{Dataset Description}
The dataset is organized into four table types:

\subsubsection{Market Candles Tables}
These tables represent historical candlestick data aggregated across multiple time intervals for all securities and trading boards. The data is stored in a hierarchical folder structure, with separate subfolders for each interval: \texttt{10min}, \texttt{1hour}, \texttt{1day}, \texttt{1week}, \texttt{1month}, and \texttt{1quarter}. Within each interval folder, individual CSV files follow the naming pattern \texttt{\{ticker\}-\{interval\}.csv}. The CSV schema includes the candle start time (\texttt{begin}), the four standard price points (\texttt{open}, \texttt{high}, \texttt{low}, \texttt{close}), the total trading value in rubles (\texttt{value}), and the volume in lots (\texttt{volume}).

\subsubsection{Dividends Table (\texttt{dividends.csv})}
This table captures the historical dividend payments for securities in the dataset. Each record contains the ticker (\texttt{secid}), the international identifier (\texttt{isin}), the critical register closure date (\texttt{registryclosedate}), the dividend amount per share (\texttt{value}), and the currency of payment (\texttt{currencyid}).

\subsubsection{Market History Table (\texttt{market\_history.csv})}
This table contains daily trading records for all securities across all trading boards. Each row is uniquely identified by the composite key of ticker (\texttt{SECID}), trading board identifier (\texttt{BOARDID}), and trading date (\texttt{TRADEDATE}). The CSV schema includes opening price (\texttt{OPEN}), intraday peak (\texttt{HIGH}), intraday trough (\texttt{LOW}), and closing price (\texttt{CLOSE}), along with trading volume in lots (\texttt{VOLUME}) and total monetary value of all trades (\texttt{VALUE}). Additional metrics include the volume-weighted average price (\texttt{WAPRICE}), the number of individual transactions (\texttt{NUMTRADES}), the denomination currency (\texttt{CURRENCYID}), security's short name (\texttt{SHORTNAME}), official settlement closing price (\texttt{LEGALCLOSEPRICE}), trend closing price (\texttt{TRENDCLSPR}), market prices (\texttt{MARKETPRICE2}, \texttt{MARKETPRICE3}), admitted quote (\texttt{ADMITTEDQUOTE}), associated trade values (\texttt{MP2VALTRD}, \texttt{MARKETPRICE3TRADESVALUE}, \texttt{ADMITTEDVALUE}), weighted average value (\texttt{WAVAL}), trading session identifier (\texttt{TRADINGSESSION}), and trade session date (\texttt{TRADE\_SESSION\_DATE}).

\subsubsection{Security Descriptions Table (\texttt{security\_descriptions.csv})}
This table acts as a static reference, containing one row per ticker with unchanging metadata about each security. Key columns include the ticker (\texttt{SECID}), the issuer's full name (\texttt{NAME}), the trading short name (\texttt{SHORTNAME}), the international identifier (\texttt{ISIN}), the state registration number (\texttt{REGNUMBER}), the nominal value (\texttt{FACEVALUE}) in the specified unit (\texttt{FACEUNIT}), the total issue size (\texttt{ISSUESIZE}), the issue date (\texttt{ISSUEDATE}), the listing level (\texttt{LISTLEVEL}), and classification fields (\texttt{GROUP}, \texttt{TYPE}, \texttt{TYPENAME}). Additional metadata covers the emitter identifier (\texttt{EMITTER\_ID}), the registry date (\texttt{REGISTRY\_DATE}), trading session flags (\texttt{MORNINGSESSION}, \texttt{EVENINGSESSION}, \texttt{WEEKENDSESSION}), qualified investor status (\texttt{ISQUALIFIEDINVESTORS}, \texttt{QUALINVESTORGROUP}), high-risk flag (\texttt{HIGHRISK}), and depositary organization details (\texttt{DEPOSITARY\_ORG\_NAME}, \texttt{DEPOSITARYNOTETYPE}).

\end{document}
